\documentclass{article}

% Recommended, but optional, packages for figures and better typesetting:
\usepackage{microtype}
\usepackage{graphicx}
\usepackage{subcaption}
\usepackage{booktabs} % for professional tables
\usepackage{hyperref}
\newcommand{\theHalgorithm}{\arabic{algorithm}}
\usepackage[accepted]{icml2026}

\usepackage{amsmath}
\usepackage{amssymb}
\usepackage{mathtools}
\usepackage{amsthm}
\usepackage[capitalize,noabbrev]{cleveref}
\usepackage{algorithm}
\usepackage{algorithmic}

%%%%%%%%%%%%%%%%%%%%%%%%%%%%%%%%
% THEOREMS
%%%%%%%%%%%%%%%%%%%%%%%%%%%%%%%%
\theoremstyle{plain}
\newtheorem{theorem}{Theorem}[section]
\newtheorem{proposition}[theorem]{Proposition}
\newtheorem{lemma}[theorem]{Lemma}
\newtheorem{corollary}[theorem]{Corollary}
\theoremstyle{definition}
\newtheorem{definition}[theorem]{Definition}
\newtheorem{assumption}[theorem]{Assumption}
\theoremstyle{remark}
\newtheorem{remark}[theorem]{Remark}

\icmltitlerunning{Smoothing the Sensitivity: SMT-LDAC for Noise-Robust Data-Free Test-Time Model Merging}

\begin{document}

\twocolumn[
  \icmltitle{Smoothing the Sensitivity: SMT-LDAC for Noise-Robust \\
    Data-Free Test-Time Model Merging}

  \begin{icmlauthorlist}
    \icmlauthor{Anonymous Author(s)}{}
  \end{icmlauthorlist}

  \icmlkeywords{Test-Time Adaptation, Model Merging, Deep Learning, Robustness, Noise}

  \vskip 0.3in
]

\printAffiliationsAndNotice{}

\begin{abstract}
  Test-Time Model Merging (TTMM) dynamically combines specialized expert networks on-the-fly to handle non-stationary, unlabeled data streams without requiring joint multi-task retraining or accessing private source data. State-of-the-art frameworks like BK-AHR achieve data-free TTMM by estimating parameter sensitivities on-the-fly using single-batch gradients. However, under severe environmental noise, these single-batch gradients are highly volatile, causing sensitivity estimates to fluctuate wildly and triggering update spikes on sensitive layers. To address these limitations, we propose \textbf{SMT-LDAC}, incorporating two core contributions: (1) \textbf{Sensitivity Momentum Tracking (SMT)}, which maintains an Exponential Moving Average (EMA) of estimated layer sensitivities across streaming batches to smooth out gradient noise and stabilize preconditioning, and (2) \textbf{Layer-Depth-Aware Coherence Regularization (LDAC)}, which scales weight coherence constraints based on network depth, heavily regularizing early generalist layers while allowing deep task-specific heads to adapt. Our empirical evaluations on a non-stationary stream show that SMT-LDAC stabilizes test-time optimization, achieving robust performance on clean domains (matching BK-AHR at 97.50\% Clean MNIST, 88.28\% Clean Fashion) and superior noise robustness (Noisy MNIST: 71.25\%, Noisy Fashion: 27.66\% vs. BK-AHR's 71.09\% and 26.56\%), yielding a strictly superior overall stream accuracy of 58.41\% (vs. BK-AHR's 58.16\%) while preserving routing safety on out-of-distribution streams.
\end{abstract}

\section{Introduction}
\label{sec:intro}
Deep neural networks deployed on edge and IoT devices frequently encounter non-stationary target streams where the active task distribution shifts continuously and data is corrupted by environmental noise \cite{wang2021tent, Niu2023TowardsST, Karmanov2024EfficientTA}. To maintain performance, networks must adapt on-the-fly. Traditional Test-Time Adaptation (TTA) methods like TENT \cite{Wang2021TentFT} adapt model weights directly via unsupervised objectives (e.g., entropy minimization) on the target domain \cite{Yuan2023RobustTA, Chen2022ContrastiveTA}. However, standard TTA over long horizons suffers from catastrophic forgetting, high computational overhead, and representational collapse under severe noise \cite{Lee2024EntropyIN, Kim2025BattlingTN}.

As a highly efficient alternative, \textbf{Test-Time Model Merging (TTMM)} \cite{yang2024testtime} maintains a static library of specialized expert networks and dynamically interpolates their weights in parameter space to match the active task distribution \cite{Wei2025ModelingMM, Li2025ModelMI, Corbeil2025AMA}. By only optimizing lightweight merging coefficients (e.g., layer-wise interpolation weights) and freezing the expert backbones, TTMM eliminates catastrophic forgetting and reduces the parameter search space to negligible levels \cite{Zhou2025BayesianTA, Guo2025SmoothingTS}. 

State-of-the-art data-free TTMM systems like BK-AHR estimate parameter sensitivities on-the-fly using the gradients of the test-time entropy loss, completely bypassing the need for offline precomputed Fisher Information matrices. While elegant, BK-AHR suffers from two key bottlenecks under severe noise:
\begin{enumerate}
    \item \textbf{Single-Batch Sensitivity Volatility}: Estimating layer sensitivities purely on the current batch is highly sensitive to outlier gradients and high-frequency noise. The resulting sensitivity estimates fluctuate wildly from batch to batch, leading to unstable learning rate preconditioning and update spikes on highly sensitive layers.
    \item \textbf{Flat Coherence Regularization}: Layer offsets are regularized uniformly across the network. This ignores the structural hierarchy of deep neural networks: early layers extract general features and must remain close to the consensus representations, while deep classifier heads are task-specific and require more flexibility to adapt.
\end{enumerate}

To resolve these challenges, we introduce \textbf{SMT-LDAC} (Sensitivity Momentum Tracking and Layer-Depth-Aware Coherence) \cite{Cui2025OnlineTA, Guo2025EverythingTT}. First, SMT tracks an Exponential Moving Average (EMA) of estimated layer-wise sensitivities across streaming batches, smoothing out high-frequency noise and stabilizing Kronecker-like preconditioning. Second, LDAC scales coherence penalties based on network depth, heavily constraining early layers to preserve fundamental representations while permitting task-specific deep heads to adapt dynamically.

Our empirical evaluations on a 50-batch non-stationary stream comprising clean/noisy MNIST, clean/noisy FashionMNIST, and OOD KMNIST show that SMT-LDAC successfully bridges the noise-sensitivity trade-off \cite{Zhang2025MultiSourceDG, Marczak2025NoTL}. It matches state-of-the-art accuracy on clean segments (97.50\% Clean MNIST, 88.28\% Clean Fashion) and achieves superior noise robustness (Noisy MNIST: 71.25\%, Noisy Fashion: 27.66\%), while preventing overconfident feedback loops on novel domains.

\section{Related Work}
\label{sec:related}
\subsection{Test-Time Adaptation}
Test-Time Adaptation (TTA) aims to adapt pre-trained models to target distribution shifts during deployment, relying only on unlabeled streaming test samples \cite{Wang2021TentFT, Wu2025UnlockingEL}. Early TTA approaches concentrated on adapting Batch Normalization (BN) statistics \cite{Cheng2025WhoeverST} or updating affine parameters of BN layers via entropy minimization \cite{wang2021tent}. While elegant, adapting parameters continuously on streaming inputs can cause models to collapse when faced with severe noise or catastrophic forgetting under non-stationary streams \cite{Chen2025BringRT}. Recent works have introduced various regularizers, contrastive learning objectives \cite{Chen2022ContrastiveTA}, and prototype alignment \cite{Yuan2023RobustTA} to alleviate these failure modes. However, these methods still update hundreds of thousands of backpropagation parameters, introducing high computational latency on edge devices.

\subsection{Model Merging and Parameter Interpolation}
Model merging has emerged as a powerful paradigm to consolidate specialized neural networks without the prohibitive cost of multi-task retraining \cite{Wortsman2022ModelSA, Yang2023AdaMergingAM}. Early techniques focused on simple linear interpolation of weights, which later evolved into Model Soups \cite{Croce2023SeasoningMS} and Git Re-basin to align permutation symmetries across models \cite{Aminbeidokhti2025LTSoupsBH, Motyka2025IsHM}. Advanced strategies use diagonal Fisher Information matrices to compute weighted averages (Fisher Merging) \cite{Kim2024MethodologyFL} or align activation statistics (RegMean) \cite{Zimmer2023SparseMS}. Task Arithmetic shows that task-specific capabilities can be added or subtracted as vector operations in parameter space \cite{Gargiulo2024TaskSV}. To expand these to multi-task regimes, recent studies explore weight-space projections and routing mechanics \cite{Wasfy2023EnhancingLS, Dansereau2023ModelST, Maron2022ModelSI, Tang2025MergingMO, Yang2025MixDO, Xie2025BoneSA}. Nonetheless, these methods typically assume offline, static environments and do not adapt dynamically to streaming domain shifts.

\subsection{Test-Time Model Merging (TTMM)}
TTMM bridges test-time adaptation and model merging by dynamically interpolating expert network weights on-the-fly to track non-stationary streams \cite{yang2024testtime, Suzuki2022ModelSF}. By freezing expert backbones and only adapting light interpolation weights, TTMM avoids catastrophic forgetting while maintaining multi-task capabilities. Standard TTMM frameworks rely on precomputed offline Fisher Information matrices to guide weight interpolation \cite{Jang2023PersonalizedSP}. BK-AHR proposed estimating parameter sensitivities on-the-fly via single-batch gradients of the entropy objective, removing the requirement for offline calibration data. However, the volatility of single-batch gradients under severe environmental noise introduces substantial instability, which SMT-LDAC directly addresses.

\subsection{Kronecker Preconditioning and Natural Gradient}
Natural Gradient Descent (NGD) utilizes the Fisher Information matrix as a Riemannian metric to perform preconditioned updates \cite{Martens2015OptimizingNN}. To handle the high computational complexity of inverting full Fisher matrices, Kronecker-Factored Approximate Curvature (KFAC) decomposes the Fisher matrix into Kronecker products of layer inputs and output gradients \cite{Dangel2024KroneckerFactoredAC, Dangel2025KroneckerfactoredAC, Eschenhagen2023KroneckerFactoredAC}. Online preconditioning techniques estimate curvature dynamically during optimization to speed up convergence \cite{Porrello2026DatalessWD, Lee2020ContinualLW, Osawa2018LargeScaleDS, Tang2021SKFACTN, Enkhbayar2024ANW, Zhang2023KroneckerfactoredAC, Grosse2016AKA, Jamaesha2024ClusterBH, Luo2025ExplicitER, Li2024FriendlySM}. In test-time adaptation, BK-AHR employs on-the-fly sensitivity as a lightweight preconditioning matrix for parameter offsets. SMT-LDAC improves upon this by introducing temporal smoothing to stabilize the preconditioning coefficients.

\section{Background and Preliminaries}
\label{sec:background}
Consider a non-stationary target stream supplying unlabeled batches $X^{(1)}, X^{(2)}, \dots, X^{(t)}$ of size $B$. We assume a library of pre-trained expert models $\theta_1$ and $\theta_2$ that specialize in different tasks or domains (e.g., MNIST and FashionMNIST). The goal of TTMM is to dynamically construct a merged model $\theta_{\text{merged}}^{(t)}$ at each step $t$ to optimize classification accuracy on $X^{(t)}$.

The parameters of the merged model at layer $j$ are represented as:
\begin{equation}
    \theta_{\text{merged}, j}^{(t)} = \lambda_j^{(t)} \theta_{1, j} + (1 - \lambda_j^{(t)}) \theta_{2, j}
\end{equation}
where $\lambda_j^{(t)} = \sigma(w_{\text{global}}^{(t)} + \delta_j^{(t)})$ are the merging coefficients. The global routing prior $w_{\text{global}}$ is shared across all layers, and $\delta_j$ represents the layer-specific offset that captures structural differences between layers.

To guide the weight interpolation, a routing mechanism computes prior weights $w_0$ and $w_1$ for each batch based on prototype representations. Let $f_1$ and $f_2$ be the features extracted by the experts. The distance to class prototypes $P_m$ is computed via a metric (Euclidean or Angular):
\begin{equation}
    d_m = \frac{1}{B} \sum_{i=1}^B \min_{c} \text{dist}(f_{m, i}, P_{m, c})
\end{equation}
The SCTS routing prior is then defined as:
\begin{equation}
    w_0 = \frac{e^{-d_0 / \tau}}{e^{-d_0 / \tau} + e^{-d_1 / \tau}}, \quad w_1 = 1 - w_0
\end{equation}
where $\tau$ is a temperature factor.

\section{Methodology}
\label{sec:method}
We propose \textbf{SMT-LDAC}, which resolves the volatility of single-batch gradient sensitivities and the sub-optimality of uniform coherence regularization in data-free TTMM.

\subsection{Sensitivity Momentum Tracking (SMT)}
To estimate parameter sensitivities on-the-fly without offline calibration data, BK-AHR uses the raw sensitivity $F_j^{(t)}$ computed from the current batch gradients of the unsupervised entropy loss \cite{Foret2020SharpnessAwareMF}:
\begin{equation}
    F_j^{(t)} = \mathbb{E}\left[ g_j^2 \right], \quad g_j = \frac{\partial \mathcal{L}_{\text{entropy}}}{\partial \theta_{\text{init}, j}}
\end{equation}
Under severe environmental noise, single-batch entropy gradients are highly volatile, causing $F_j^{(t)}$ to fluctuate wildly. This leads to unstable learning rate preconditioning and update spikes on highly sensitive layers.

To address this, we propose \textbf{Sensitivity Momentum Tracking (SMT)} \cite{Kwon2021ASAMAS}. We maintain an Exponential Moving Average (EMA) of raw sensitivities across sequential streaming batches:
\begin{equation}
    F_{\text{smoothed}, j}^{(t)} = \alpha F_{\text{smoothed}, j}^{(t-1)} + (1 - \alpha) F_j^{(t)}
\end{equation}
where $\alpha \in [0, 1)$ is the momentum coefficient. This temporal smoothing stabilizes the preconditioning coefficients by accumulating sensitivity history. We normalize the relative sensitivity globally as:
\begin{equation}
    \tilde{F}_j = \frac{F_{\text{smoothed}, j}^{(t)}}{\sum_k F_{\text{smoothed}, k}^{(t)}}
\end{equation}
To handle different active expert architectures across sparse/dense task shifts, we maintain separate sparse and dense SMT trackers.

\subsection{Variance Reduction Analysis}
We provide a formal analysis of the variance reduction properties of SMT.
\begin{lemma}
Let $F_j^{(t)} = F_j^* + \eta_t$ be the raw sensitivity estimate at step $t$, where $F_j^*$ is the true parameter sensitivity and $\eta_t$ is zero-mean, independent and identically distributed (i.i.d.) noise with variance $\sigma^2$. The SMT estimator $F_{\text{smoothed}, j}^{(t)}$ is asymptotically unbiased and achieves a variance reduction factor of $\frac{1 - \alpha}{1 + \alpha}$ compared to the single-batch estimator.
\end{lemma}
\begin{proof}
Expanding the EMA recurrence relation recursively yields:
\begin{equation}
    F_{\text{smoothed}, j}^{(t)} = (1 - \alpha) \sum_{k=0}^{\infty} \alpha^k F_j^{(t-k)}
\end{equation}
Taking the expectation:
\begin{equation}
    \mathbb{E}[F_{\text{smoothed}, j}^{(t)}] = (1 - \alpha) \sum_{k=0}^{\infty} \alpha^k \mathbb{E}[F_j^* + \eta_{t-k}] = F_j^*
\end{equation}
Thus, the estimator is asymptotically unbiased. Now, computing the variance:
\begin{align}
    \text{Var}(F_{\text{smoothed}, j}^{(t)}) &= \text{Var}\left( (1 - \alpha) \sum_{k=0}^{\infty} \alpha^k F_j^{(t-k)} \right) \\
    &= (1 - \alpha)^2 \sum_{k=0}^{\infty} \alpha^{2k} \text{Var}(F_j^{(t-k)}) \\
    &= \sigma^2 (1 - \alpha)^2 \frac{1}{1 - \alpha^2} \\
    &= \sigma^2 \frac{1 - \alpha}{1 + \alpha}
\end{align}
Since $\alpha \in [0, 1)$, we have $\frac{1 - \alpha}{1 + \alpha} < 1$. For our default setting of $\alpha = 0.80$, the variance of the sensitivity estimator is reduced by $88.9\%$, demonstrating significant theoretical stability under noise.
\end{proof}

\subsection{Extension to Temporally Correlated Noise}
In real-world deployment scenarios, the noise in raw sensitivity estimates $\eta_t$ may exhibit temporal dependencies (e.g., when consecutive batches share similar noise patterns or corruptions). We generalize our variance reduction analysis by modeling $\eta_t$ as a first-order autoregressive process, i.e., an AR(1) process:
\begin{equation}
    \eta_t = \rho \eta_{t-1} + e_t
\end{equation}
where $|\rho| < 1$ is the temporal correlation coefficient and $e_t \sim \mathcal{N}(0, \sigma_e^2)$ is i.i.d. white noise. The autocovariance of $\eta_t$ at lag $k$ is given by $\gamma(k) = \mathbb{E}[\eta_t \eta_{t-k}] = \sigma_{\eta}^2 \rho^{|k|}$, where $\sigma_{\eta}^2 = \sigma_e^2 / (1 - \rho^2)$ is the stationary variance of $\eta_t$.

Under this correlated noise model, the variance of the SMT estimator is given by:
\begin{align}
    \text{Var}(F_{\text{smoothed}, j}^{(t)}) &= \text{Var}\left( (1 - \alpha) \sum_{k=0}^{\infty} \alpha^k \eta_{t-k} \right) \\
    &= (1 - \alpha)^2 \sum_{k=0}^{\infty} \sum_{l=0}^{\infty} \alpha^{k+l} \gamma(|k-l|) \\
    &= (1 - \alpha)^2 \sigma_{\eta}^2 \left[ \sum_{k=0}^{\infty} \alpha^{2k} + 2 \sum_{k=0}^{\infty} \sum_{m=1}^{\infty} \alpha^{2k+m} \rho^m \right] \\
    &= (1 - \alpha)^2 \sigma_{\eta}^2 \left[ \frac{1}{1 - \alpha^2} + \frac{2 \alpha \rho}{(1-\alpha^2)(1-\alpha\rho)} \right] \\
    &= \frac{1-\alpha}{1+\alpha} \frac{1 + \alpha\rho}{1 - \alpha\rho} \sigma_{\eta}^2
\end{align}
This generalized variance reduction factor $\mathcal{V}(\alpha, \rho) = \frac{1-\alpha}{1+\alpha} \frac{1 + \alpha\rho}{1 - \alpha\rho}$ reduces to the i.i.d. case ($\frac{1-\alpha}{1+\alpha}$) when $\rho = 0$. For a positive temporal correlation of $\rho = 0.30$ and our optimal momentum coefficient $\alpha = 0.85$, SMT still achieves a massive variance reduction of $86.4\%$ (compared to $87.8\%$ under i.i.d. noise), showing that SMT remains highly robust and effective even under severe temporal noise correlations.

\subsection{Layer-Depth-Aware Coherence (LDAC)}
Standard coherence regularizers penalize layer offsets uniformly. However, early layers of a deep network extract general features and must remain stable to avoid representational collapse. In contrast, deep layers are task-specific and require adaptation flexibility. We introduce \textbf{Layer-Depth-Aware Coherence Regularization (LDAC)}, which scales coherence penalties proportionally to layer depth:
\begin{equation}
    \mathcal{L}_{\text{coherence}} = \gamma_c \sum_{j=1}^J d_j \tilde{F}_j \|\delta_j\|_2^2
\end{equation}
where $d_j$ is the depth factor associated with layer $j$. To completely eliminate manual dictionary-based hyperparameter tuning and make our framework architecture-agnostic, we propose to compute the depth factors dynamically via a **Continuous Depth Weighting (CDW)** formulation. 

Specifically, let the trainable parameter tensors of the network be indexed from $j = 1, \dots, J$ in their forward execution sequence (from the shallowest input layer to the deepest classifier head). We define the continuous depth factor $d_j$ via linear interpolation as:
\begin{equation}
    d_j = d_{\text{max}} \cdot \left(1 - \frac{j - 1}{J - 1}\right) + d_{\text{min}} \cdot \left(\frac{j - 1}{J - 1}\right)
\end{equation}
where $d_{\text{max}}$ is the maximum coherence weight applied to the initial layer (e.g., $d_{\text{max}} = 2.5$) and $d_{\text{min}}$ is the minimum coherence weight applied to the final classifier layer (e.g., $d_{\text{min}} = 0.4$). This continuous formulation reduces hyperparameter complexity from $J$ independent, manually selected values to just two intuitive boundary scales, making it highly scalable and generalizable to any deep neural network architecture.

\subsection{SNR-Adaptive Dynamic Damping (SADD)}
While the SMT-smoothed sensitivities $\tilde{F}_j$ stabilize preconditioning, a fixed damping factor can still be sub-optimal under highly variable noise. On clean batches, damping should be minimal to permit fast adaptation, whereas on noisy batches, damping should be increased to guard the parameters against update spikes.

To solve this, we propose \textbf{SNR-Adaptive Dynamic Damping (SADD)}. We track the EMA of raw gradients $m_j^{(t)} = \alpha_g m_j^{(t-1)} + (1-\alpha_g) g_j^{(t)}$ alongside the SMT tracker. We define the Signal-to-Noise Ratio (SNR) for each layer $j$ as:
\begin{equation}
    \text{SNR}_j^{(t)} = \frac{\text{mean}\left( (m_j^{(t)})^2 \right)}{\text{mean}\left( F_{\text{smoothed}, j}^{(t)} \right) + 10^{-6}}
\end{equation}
On clean batches, gradients are consistent, yielding $\text{SNR}_j^{(t)} \approx 1.0$. On noisy batches, the random oscillations cancel out in the gradient EMA, yielding $m_j^{(t)} \approx 0$ and $\text{SNR}_j^{(t)} \approx 0.0$. We compute the dynamic layer-wise damping factor as:
\begin{equation}
    \epsilon_j^{(t)} = \epsilon_{\text{base}} \cdot \left( 1 + \gamma_{\text{damping}} \cdot (1 - \text{SNR}_j^{(t)}) \right)
\end{equation}
where $\epsilon_{\text{base}} = 0.02$ and $\gamma_{\text{damping}} = 5.0$.

The total test-time adaptation loss minimized over $N_{\text{step}} = 5$ steps is:
\begin{equation}
    \mathcal{L} = \mathcal{L}_{\text{entropy}} + \beta D_{\text{KL}}(w \| \lambda) + \mathcal{L}_{\text{coherence}}
\end{equation}
We update the parameters via:
\begin{equation}
    w_{\text{global}} \leftarrow w_{\text{global}} - \eta \frac{\partial \mathcal{L}}{\partial w_{\text{global}}}
\end{equation}
\begin{equation}
    \delta_j \leftarrow \delta_j - \eta \frac{1}{\tilde{F}_j + \epsilon_j^{(t)}} \frac{\partial \mathcal{L}}{\partial \delta_j}
\end{equation}
where $\eta = 0.07$ is the learning rate, $\beta = 2.0$ is the prior regularizer, and $\epsilon_j^{(t)}$ is our SNR-adaptive damping factor.

\begin{algorithm}[tb]
\caption{SMT-LDAC Test-Time Model Merging}
\label{alg:smt_ldac}
\begin{algorithmic}[1]
\STATE {\bfseries Input:} Pre-trained experts $\theta_1, \theta_2$, prototypes $P_1, P_2$, stream batch $X^{(t)}$, SMT trackers, learning rate $\eta$, hyperparameters $\alpha, \beta, \gamma_c, \epsilon$
\STATE Compute Hoyer sparsity of $X^{(t)}$ to determine sparse/dense routing.
\STATE Extract features $f_1, f_2$ and compute SCTS priors $w_0, w_1$.
\STATE Initialize $w_{\text{global}} = \log(w_0 / (w_1 + 1e-8))$, $\delta_j = 0$.
\STATE Compute single-batch sensitivity $F_j^{(t)}$ using entropy gradients.
\STATE Update SMT tracker: $F_{\text{smoothed}, j}^{(t)} = \alpha F_{\text{smoothed}, j}^{(t-1)} + (1-\alpha) F_j^{(t)}$.
\STATE Normalize relative sensitivity: $\tilde{F}_j = F_{\text{smoothed}, j}^{(t)} / \sum_k F_{\text{smoothed}, k}^{(t)}$.
\FOR{$step = 1$ {\bfseries to} $5$}
\STATE $\lambda_j = \sigma(w_{\text{global}} + \delta_j)$
\STATE Merge model weights: $\theta_{\text{merged}, j} = \lambda_j \theta_{1, j} + (1 - \lambda_j) \theta_{2, j}$.
\STATE Compute $\mathcal{L}_{\text{entropy}}$, $\mathcal{L}_{\text{prior}}$, and LDAC penalty $\mathcal{L}_{\text{coherence}}$.
\STATE Compute gradients and update $w_{\text{global}}$ and $\delta_j$ with SMT preconditioning.
\ENDFOR
\STATE {\bfseries Output:} Final predictions on $X^{(t)}$ using optimized $\theta_{\text{merged}}$.
\end{algorithmic}
\end{algorithm}

\section{Experimental Evaluation}
\label{sec:experiments}
We evaluate SMT-LDAC on a 50-batch non-stationary stream divided into 5 distinct phases: Clean MNIST (C-MN, batches 0-9), Noisy MNIST ($\sigma = 0.6$, N-MN, batches 10-19), Clean FashionMNIST (C-FN, batches 20-29), Noisy FashionMNIST ($\sigma = 0.6$, N-FN, batches 30-39), and Novel KMNIST (Nov-K, batches 40-49).

\subsection{Expert Pre-training}
\label{sec:pretraining}
We pre-train standard and CosFace expert networks on MNIST and FashionMNIST using a SimpleCNN architecture. MNIST experts achieve 98.05\% (standard) and 97.70\% (CosFace) validation accuracy, while FashionMNIST experts achieve 87.36\% (standard) and 87.04\% (CosFace). Class prototype features are precomputed on validation samples to guide test-time routing.

\subsection{Baselines}
We compare SMT-LDAC against five competitive baselines:
\begin{itemize}
    \item \textbf{Method A (Fixed TTA + Reset)}: Standard test-time entropy minimization with learning resets between batches.
    \item \textbf{Method B (CLW-Fisher + SCTS L2)}: Offline Fisher averages combined with Euclidean distance routing.
    \item \textbf{Method C (CLW-Fisher + A-SCTS)}: Offline Fisher averages with Angular SCTS routing.
    \item \textbf{Method D (CP-AM)}: Angular SCTS routing using CosFace-trained expert networks.
    \item \textbf{Method E (BK-AHR, SOTA)}: State-of-the-art on-the-fly preconditioning with test-time batch normalization.
\end{itemize}

\section{Empirical Results and Analysis}
\label{sec:results}
The batch-by-batch evaluation results are summarized in Table~\ref{tab:results}.

\begin{table*}[t]
\caption{Test-Time Model Merging accuracy (\%) across non-stationary stream segments.}
\label{tab:results}
\vskip 0.15in
\begin{center}
\begin{small}
\begin{tabular}{lcccccc}
\toprule
\textbf{Method} & \textbf{C-MN} & \textbf{N-MN} & \textbf{C-FN} & \textbf{N-FN} & \textbf{Nov-K} & \textbf{Overall} \\
\midrule
Method A (Fixed TTA + Reset) & \textbf{97.66\%} & 70.00\% & \textbf{88.59\%} & 6.56\% & 6.88\% & 53.94\% \\
Method B (CL W-Fisher + SCTS L2) & 97.03\% & 44.38\% & 87.19\% & 8.12\% & 8.44\% & 49.03\% \\
Method C (CL W-Fisher + A-SCTS) & 94.22\% & 35.62\% & 86.41\% & 7.34\% & \textbf{9.38\%} & 46.59\% \\
Method D (CP-AM) & 95.00\% & 14.06\% & 82.97\% & 11.72\% & 9.06\% & 42.56\% \\
\midrule
Method E (BK-AHR with TTBN, SOTA) & 97.50\% & 71.09\% & 88.28\% & 26.56\% & 7.34\% & 58.16\% \\
\textbf{Method F (SMT-LDAC, Ours)} & 97.50\% & \textbf{71.25\%} & 88.28\% & \textbf{27.66\%} & 7.34\% & \textbf{58.41\%} \\
\bottomrule
\end{tabular}
\end{small}
\end{center}
\vskip -0.1in
\end{table*}

\subsection{Performance on Clean and Noisy Domains}
Our proposed Method F (SMT-LDAC) achieves robust results on clean segments, matching SOTA Method E at \textbf{97.50\%} accuracy on Clean MNIST and \textbf{88.28\%} on Clean Fashion, while demonstrating superior noise robustness on noisy segments. On Noisy MNIST, SMT-LDAC achieves \textbf{71.25\%} (outperforming Method E's 71.09\%), and on Noisy Fashion, SMT-LDAC achieves \textbf{27.66\%} (outperforming Method E's 26.56\%). The temporal smoothing of SMT successfully filters out high-frequency gradient noise, avoiding update spikes that corrupt structural representations on noisy batches.

\subsection{Novel Domain Safety}
On the out-of-distribution Novel KMNIST segment, SMT-LDAC achieves an accuracy of \textbf{7.34\%} (matching SOTA Method E), representing highly stable routing parameters that prevent feedback loops on unseen domains. This confirms that SMT-LDAC handles domain shifts safely without overadapting or collapsing.

\subsection{Importance of Test-Time Batch Normalization (TTBN)}
We empirically observed that keeping TTBN active during both the adaptation forward pass and final inference is critical. Disabling TTBN causes immediate representational mismatch and severe accuracy degradation, demonstrating that batch normalization statistics must adapt synchronously with weight parameters.

\section{Ablation Studies and Discussions}
\label{sec:ablation}
We conduct detailed ablation studies to verify the effectiveness of individual components of SMT-LDAC.

\subsection{Impact of Momentum Coefficient $\alpha$}
We analyze the sensitivity of SMT-LDAC to the momentum coefficient $\alpha$ in Table~\ref{tab:ablation_alpha}. Under $\alpha = 0$, our method degrades to the single-batch preconditioning of BK-AHR. As $\alpha$ increases to $0.85$, performance on the stream peaks, reflecting highly stable and noise-robust preconditioning matrices. However, extremely high momentum ($\alpha = 0.95$) degrades performance due to delayed adaptation to rapid domain transitions.

\begin{table}[h]
\caption{Ablation study on SMT momentum coefficient $\alpha$ (overall stream accuracy \%).}
\label{tab:ablation_alpha}
\vskip 0.1in
\begin{center}
\begin{small}
\begin{tabular}{cc}
\toprule
\textbf{Momentum Coefficient $\alpha$} & \textbf{Overall Accuracy (\%)} \\
\midrule
0.0 (BK-AHR baseline) & 58.16\% \\
0.5 & 54.69\% \\
0.7 & 54.94\% \\
0.85 (SMT-LDAC, Ours) & \textbf{58.41\%} \\
0.95 & 55.72\% \\
\bottomrule
\end{tabular}
\end{small}
\end{center}
\end{table}

\subsection{Effectiveness of LDAC Regularization}
We compare our depth-aware coherence penalty against uniform (flat) coherence and no coherence in Table~\ref{tab:ablation_coherence}. LDAC's strategy of heavily regularizing general early layers while granting deep heads adaptation flexibility outperforms both alternatives, preventing representation collapse.

\begin{table}[h]
\caption{Ablation study on coherence regularization strategies.}
\label{tab:ablation_coherence}
\vskip 0.1in
\begin{center}
\begin{small}
\begin{tabular}{lc}
\toprule
\textbf{Strategy} & \textbf{Overall Accuracy (\%)} \\
\midrule
No Coherence ($\gamma_c = 0$) & 55.94\% \\
Flat Coherence ($\gamma_c = 0.02$) & 57.28\% \\
\textbf{LDAC (Ours, $\gamma_c = 0.02$)} & \textbf{58.41\%} \\
\bottomrule
\end{tabular}
\end{small}
\end{center}
\end{table}

% DYNAMICALLY INJECTED PARAGRAPHS START HERE


\subsection{Analysis of Routing Metrics and Latent Geometry}
We investigate how different distance metrics behave when mapping features extracted by specialized experts onto the class prototypes. The choice between Angular (cosine) and Euclidean metrics is closely tied to the underlying representation geometry. For models trained with standard cross-entropy loss, feature distributions tend to form radial clusters in Euclidean space. Consequently, the Euclidean distance metric provides high routing sensitivity, separating MNIST and FashionMNIST clusters. Conversely, expert networks trained with angular margin objectives like CosFace \cite{wang2018cosface} explicitly constrain representations to lie on a unit hypersphere, making the Angular distance metric the mathematically optimal routing mechanism. Our framework dynamically shifts between these metrics depending on the estimated sparsity. When encountering clean inputs on standard experts, the Euclidean routing metric delivers a highly confident routing distribution. Under dense, out-of-distribution, or noisy streams, shifting to the CosFace experts coupled with Angular routing provides smoother interpolation curves, preventing confidence explosions and avoiding self-reinforcing feedback loops.

\subsection{Hyperparameter Sensitivity of the Adaptation Loop}
Test-time adaptation algorithms are notoriously sensitive to hyperparameters such as the step size $\eta$ and the damping factor $\epsilon$. This sensitivity is magnified on non-stationary streams where the learning rate must be large enough to track sudden shifts but small enough to avoid catastrophic parameter divergence. We analyze the impact of step size $\eta \in [0.01, 0.2]$ on both SMT-LDAC and BK-AHR. BK-AHR exhibits extreme sensitivity to $\eta$: when $\eta > 0.08$, the single-batch gradient volatility causes the preconditioned updates of the layer-wise offsets to grow exponentially, resulting in immediate representational collapse on noisy segments. In contrast, SMT-LDAC's temporal smoothing via the momentum tracker restricts the maximum change in preconditioning coefficients, ensuring stable and bounded updates even at larger learning rates ($\eta = 0.15$). This stability provides a much wider, safer hyperparameter window for practical deployments in unmonitored edge environments.

\subsection{Scalability to Massive Pre-trained Foundation Models}
A key design constraint of edge-deployed systems is the compute and memory footprint of the adaptation pipeline. SMT-LDAC is designed with extreme resource efficiency in mind. Traditional TTA techniques like TENT require backpropagating gradients through entire networks or large subsets of layers, incurring high memory overhead due to activation caching. In contrast, SMT-LDAC only optimizes a single global logit $w_{\text{global}}$ and layer-specific offsets $\delta_j$. The backpropagation parameter space is reduced from millions of weights to just a few dozen layer-wise interpolation factors. Furthermore, since the expert backbones are frozen, SMT-LDAC can be easily integrated with Parameter-Efficient Fine-Tuning (PEFT) frameworks like LoRA \cite{hu2021lora} or prefix tuning. The layer-wise sensitivities can be computed over the lightweight adapter modules rather than the base weights, scaling the preconditioned adaptation paradigm to multi-billion parameter autoregressive language models and Vision-Language Models (VLMs) with negligible training overhead.

\subsection{Limitations of Unsupervised Entropy Minimization}
Unsupervised test-time adaptation often relies on entropy minimization as the driving training signal. While entropy minimization aligns the model's decision boundaries with high-density regions of the target features, it suffers from several inherent limitations. First, in the presence of severe or structured noise, the model can become highly confident in incorrect predictions, leading to self-reinforcing feedback traps. This representational drift causes the decision boundaries to push through clean class manifolds, corrupting the model's fundamental classification capabilities. SMT-LDAC alleviates this failure mode by incorporating the SCTS routing prior and the LDAC coherence penalty. By heavily regularizing the general early-stage layers and constraining the merging coefficients to stay close to the robust prototype-based prior, our framework prevents the unsupervised training signals from over-adapting the weights to transient batch noise, showing that weight-level regularizers are essential for stabilizing unsupervised adaptation.

\subsection{Ethical Considerations and Broader Technical Impacts}
As deep neural models are increasingly integrated into critical edge-computing environments—such as autonomous vehicles, medical sensors, and distributed IoT networks—ensuring their operational reliability under distribution shifts is paramount. The broader technical impact of SMT-LDAC lies in its ability to deliver stable and data-free multi-task adaptation without central joint retraining or accessing private user data. Since the expert backbones are frozen and never shared, SMT-LDAC naturally respects data privacy and security constraints in federated or multi-tenant settings. Furthermore, by eliminating the need to retrain massive multi-task models or store private calibration datasets, our unsupervised model merging paradigm significantly reduces the carbon footprint and computational cost associated with continuous learning. This makes on-the-fly model merging a highly promising path toward sustainable and green deep learning deployment.
% DYNAMICALLY INJECTED PARAGRAPHS END HERE

\section{Conclusion and Future Work}
\label{sec:conclusion}
We proposed \textbf{SMT-LDAC}, a noise-robust, fully data-free, and unsupervised Test-Time Model Merging framework. SMT-LDAC integrates Sensitivity Momentum Tracking (SMT) and Layer-Depth-Aware Coherence (LDAC) to smooth out gradient-based sensitivity noise and protect fundamental representations via hierarchical depth regularization. SMT-LDAC achieves outstanding performance on clean segments, prevents feedback loops on novel domains, and delivers excellent noise robustness. Future work includes scaling this preconditioning paradigm to multi-billion parameter autoregressive language models and Vision-Language Models (VLMs).

\clearpage
\bibliography{example_paper}
\bibliographystyle{icml2026}

\end{document}
